\documentclass[11pt]{article}
\usepackage[margin=1in]{geometry}
\usepackage{xcolor}
\usepackage{tcolorbox}
\usepackage{hyperref}
\usepackage{enumitem}
\usepackage{booktabs}
\usepackage{array}
\usepackage{amsmath}

% Define OSU orange color
\definecolor{osaorange}{RGB}{220,115,38}

% Configure hyperref colors
\hypersetup{
    colorlinks=true,
    linkcolor=blue,
    urlcolor=blue,
    citecolor=blue
}

\title{}
\author{}
\date{}

\begin{document}

% Orange header box with black outline
\begin{tcolorbox}[colback=osaorange,colframe=black,boxrule=2pt,arc=0pt,left=6pt,right=6pt,top=10pt,bottom=10pt]
\centering
\textcolor{white}{\textbf{\Large Week 8 Assignment}}\\\
\textcolor{white}{\textbf{\large Contingency Tables and Chi-Square Tests}}
\end{tcolorbox}

\vspace{8pt}

\noindent\textbf{Course:} H524 Introduction to Biostatistics\\
\textbf{Total Points:} 100 points\\
\textbf{Submission:} Canvas quiz

\vspace{8pt}

\section*{Instructions}

This assignment has three parts:

\begin{itemize}
\item \textbf{Part A:} Multiple choice questions (70 points)
\item \textbf{Part B:} R computational problems (30 points) - You will need to write and run R code
\item \textbf{Part C:} Group project check-in (completion only, ungraded)
\end{itemize}

Submit your answers via the Canvas quiz.

\vspace{12pt}

\section{Part A: Multiple Choice (70 points)}

\subsection*{Contingency Table Basics (10 points)}
A contingency table displays the relationship between:

\begin{enumerate}[label=\Alph*.]
\item One categorical variable
\item Two continuous variables
\item Two categorical variables
\item One categorical and one continuous variable
\end{enumerate}

\subsection*{Expected Frequencies (10 points)}
In a chi-square test, the expected frequency for a cell is calculated as:

\begin{enumerate}[label=\Alph*.]
\item (Row total × Column total) / Grand total
\item (Row total + Column total) / 2
\item Grand total / Number of cells
\item Observed frequency / Total sample size
\end{enumerate}

\subsection*{Chi-Square Test Assumptions (10 points)}
Which of the following is a requirement for conducting a chi-square test?

\begin{enumerate}[label=\Alph*.]
\item All expected cell frequencies must be at least 10
\item All expected cell frequencies should be at least 5
\item The sample must be normally distributed
\item Variables must be continuous
\end{enumerate}

\subsection*{Degrees of Freedom (10 points)}
For a 3×4 contingency table (3 rows, 4 columns), what are the degrees of freedom for the chi-square test?

\begin{enumerate}[label=\Alph*.]
\item 7
\item 6
\item 12
\item 11
\end{enumerate}

\subsection*{Interpreting Chi-Square Results (10 points)}
A chi-square test yields $p = 0.032$. At $\alpha = 0.05$, what can we conclude?

\begin{enumerate}[label=\Alph*.]
\item There is no association between the variables
\item There is a significant association between the variables
\item The variables are normally distributed
\item The test assumptions are violated
\end{enumerate}

\subsection*{2×2 Table Setup (10 points)}
In a case-control study examining smoking and lung cancer, which variable should form the rows of a 2×2 table?

\begin{enumerate}[label=\Alph*.]
\item Smoking status (exposed/unexposed)
\item Lung cancer status (case/control)
\item Either variable can form rows
\item Neither - case-control studies cannot use 2×2 tables
\end{enumerate}

\subsection*{Fisher's Exact Test (10 points)}
When should Fisher's exact test be used instead of the chi-square test?

\begin{enumerate}[label=\Alph*.]
\item When sample size is very large ($n > 1000$)
\item When expected cell frequencies are small (less than 5 in one or more cells)
\item When variables are continuous
\item When there are more than two categories
\end{enumerate}

\newpage

\section{Part B: R Computational Problems (30 points)}

\textbf{Note:} Some problems will require you to use R to compute the answer.

\subsection*{Question 8: Fisher's Exact Test in R (15 points)}

A small clinical trial tests a new treatment for a rare disease with the following results:

\begin{center}
\begin{tabular}{l|cc}
& \textbf{Improved} & \textbf{Not Improved} \\
\hline
\textbf{Treatment} & 7 & 3 \\
\textbf{Control} & 2 & 8 \\
\end{tabular}
\end{center}

Use R to perform Fisher's exact test on this data. What is the two-sided p-value?

\begin{enumerate}[label=\Alph*.]
\item 0.0032
\item 0.0698
\item 0.1524
\item 0.5000
\end{enumerate}

\vspace{1cm}

\subsection*{Question 9: Chi-Square Test in R (15 points)}

A cohort study follows 500 smokers and 500 non-smokers for 10 years to assess heart disease incidence:

\begin{center}
\begin{tabular}{l|cc}
& \textbf{Heart Disease} & \textbf{No Heart Disease} \\
\hline
\textbf{Smoker} & 75 & 425 \\
\textbf{Non-smoker} & 35 & 465 \\
\end{tabular}
\end{center}

Use R to perform a chi-square test on this data (without continuity correction). What is the chi-square test statistic?

\begin{enumerate}[label=\Alph*.]
\item 8.234
\item 12.156
\item 16.343
\item 21.450
\end{enumerate}

\newpage

\section{Part C: Group Project Check-In (Completion Only, Ungraded)}

\subsection*{Progress Update Check-In}

\textbf{GROUP LEADER SUBMITS} a brief progress update answering:

\begin{enumerate}
\item \textbf{Has your group met at least once since Week 7?} (Yes/No)
\item \textbf{Have you loaded and explored your dataset in R?} (Yes/No)
\item \textbf{Have you consulted at least one AI tool about your data?} (Yes/No - if yes, which one(s)?)
\item \textbf{Any blockers or questions?} (Optional)
\end{enumerate}

\textbf{Reminder:} This Wednesday (12:00-1:20pm) is Office Hours Day - drop by classroom/office if your group needs help!

\textbf{Submit as:} Text entry (3-5 sentences).

\vspace{12pt}

\section*{Summary}

\begin{itemize}
\item Part A (Multiple Choice): 70 points
\item Part B (R Computational): 30 points
\item Part C (Group Check-In): Completion only (ungraded)
\item \textbf{Total: 100 points}
\end{itemize}

\end{document}
